%% ===================================================================
%% PART IV
%% ===================================================================
\part{Actual occurrence fibres, consequences and literature}

%% -------------------------------------------------------------------
\section{A worked example: two values, two regimes}\label{sec:example}

This section follows two explicit rational numbers through the whole theory. They are
\[
  x_0=\frac{5^{10}\,3^{10}}{11},\qquad x_1=\frac{5^{12}\,3^{10}}{11}.
\]
They are different numbers ($x_1=25\,x_0$) with the same denominator, the same active primes and the same relation lattice. The first sits in the exceptional box and has two holes; the second sits in the nearest larger box and is convex. All numerical data below are exact integers or fractions; they are reproduced by the script \texttt{scripts/example.py} (Appendix~\ref{app:comp}), and are checks of an example, not proofs of the theorems that explain it.

\subsection{\texorpdfstring{The occurrences of $x_0$}{The occurrences of x0}}

By Proposition~\ref{prop:bij}, the occurrences of $x_0$ correspond to the integers $\kappa\ge2$ with $\kappa^2\mid5^{10}3^{10}$ and $\kappa\equiv1\pmod{11}$. Such a $\kappa$ has the form $5^i3^j$ with $0\le i,j\le5$ and $(i,j)\ne(0,0)$. The multiplicative orders modulo $11$ are $\ord_{11}(5)=\ord_{11}(3)=5$ (note $5^5=3125=284\cdot11+1$ and $3^5=243=22\cdot11+1$). We have $3^3=27\equiv5\pmod{11}$, so
\[
  5^i3^j\equiv3^{3i+j}\pmod{11},
\]
which is $1$ if and only if $3i+j\equiv0\pmod5$, that is $j\equiv-3i\equiv2i\pmod5$. This is the planar normalization of Proposition~\ref{prop:norm} with
\[
  D_0=\langle3\rangle=\langle5\rangle,\quad d=5,\quad\lambda_1=\lambda_2=1,\quad g_1=5=3^3=g_2^{-s}\ \Longrightarrow\ s\equiv2\pmod5,
\]
and $t=s^{-1}\equiv3$, $a=d-s=3$, $c=d-t=2$ (the constants of Section~\ref{sec:cells}). The box is $Q=(5,5)=(d,d)$, the exceptional box, because $q_1=q_2=5=o_1=o_2$ lies on the arithmetic plateau of Corollary~\ref{cor:plateau}. Since $d=5$ and $s=2\notin\{1,4\}$, the pair $(d,s)$ is nonspecial.

The fibre is
\[
  F_{x_0}=F_{5,2,(5,5)}=\{(0,5),(1,2),(2,4),(3,1),(4,3),(5,0),(5,5)\},
\]
seven points, shown in Figure~\ref{fig:fibre}. The corresponding occurrences are listed in Table~\ref{tab:example}. For instance $(i,j)=(1,2)$ gives $\kappa=5\cdot9=45=4\cdot11+1$, $N=44$ and
\[
  M=\frac{a}{\kappa^2}\cdot\frac{\kappa-1}{b}=\frac{5^{10}3^{10}}{45^2}\cdot\frac{44}{11}=5^8\,3^6\cdot4=1\,139\,062\,500,
\]
and one checks directly that $M(N+1)^2/N=1\,139\,062\,500\cdot2025/44=5^{10}3^{10}/11$.

\begin{table}[ht]
\centering\small
\begin{tabular}{@{}cr rr rr@{}}
\toprule
$(i,j)$ & $\kappa=5^i3^j$ & $N=\kappa-1$ & $M$ for $x_0$ & $M$ for $x_1$\\
\midrule
$(0,5)$ & $243$ & $242$ & $214\,843\,750$ & $5\,371\,093\,750$\\
$(1,2)$ & $45$ & $44$ & $1\,139\,062\,500$ & $28\,476\,562\,500$\\
$(2,4)$ & $2025$ & $2024$ & $25\,875\,000$ & $646\,875\,000$\\
$(3,1)$ & $375$ & $374$ & $139\,421\,250$ & $3\,485\,531\,250$\\
$(4,3)$ & $16\,875$ & $16\,874$ & $3\,106\,350$ & $77\,658\,750$\\
$(5,0)$ & $3125$ & $3124$ & $16\,769\,916$ & $419\,247\,900$\\
$(5,5)$ & $759\,375$ & $759\,374$ & $69\,034$ & $1\,725\,850$\\
$(6,2)$ & $140\,625$ & $140\,624$ & --- & $9\,319\,536$\\
\bottomrule
\end{tabular}
\caption{The occurrences $(N,M)$ of $x_0=5^{10}3^{10}/11$ (seven rows, column ``$M$ for $x_0$'') and of $x_1=5^{12}3^{10}/11$ (all eight rows). The row $(6,2)$ appears only for $x_1$, since $\kappa^2=5^{12}3^4$ divides $5^{12}3^{10}$ but does not divide $5^{10}3^{10}$.}\label{tab:example}
\end{table}

\subsection{\texorpdfstring{Graver vectors and the realized union of $x_0$}{Graver vectors and the realized union of x0}}

The six standard vectors of Lemma~\ref{lem:std-graver} are Graver vectors of $\Gamma_{5,2}$, and an exhaustive search shows that there are no others (this is special to small moduli: for example $\Gamma_{7,3}$ has the further Graver vector $(3,2)$, Appendix~\ref{app:graver}). Up to sign they are:
\[
  (5,0),\ (0,5),\ (1,2),\ (1,-3),\ (3,1),\ (2,-1).
\]
Each has an arithmetic meaning. The relation $(1,2)$ says $5\cdot3^2=45\equiv1\pmod{11}$; the relation $(1,-3)$ says $5\equiv3^3$; the relation $(2,-1)$ says $5^2\equiv3$; and $(3,1)$ says $5^3\cdot3=375\equiv1$. These are the multiplicative relations among $5$ and $3$ modulo $11$ that cannot be split into two relations of the same sign pattern.

The hull $H_0=\conv F_{x_0}$ has vertices $(0,5),(1,2),(3,1),(5,0),(5,5)$ and area $35/2$. The lower-left chain of Lemma~\ref{lem:chain0} is $A=(0,5)$, $(1,2)$, $(3,1)$, $B=(5,0)$ with steps $(1,-3)$, $(2,-1)$, $(2,-1)$. The zonogon of Lemma~\ref{lem:zonogon} is the parallelogram with vertices $A,(1,2),B,(4,3)$ of area $10$. The additional cell of Lemma~\ref{lem:K} is $K=P\bigl((1,2);(1,2),(3,1)\bigr)$ with corners $(1,2),(2,4),(5,5),(4,3)$. The four cells of the realized union are
\begin{align*}
  C_1&=P\bigl((0,5);(1,-3),(2,-1)\bigr)\ \text{with corners }(0,5),(1,2),(3,1),(2,4),\\
  C_2&=P\bigl((1,2);(1,2),(2,-1)\bigr)\ \text{with corners }(1,2),(2,4),(4,3),(3,1),\\
  C_3&=K,\\
  C_4&=P\bigl((2,4);(1,-3),(2,-1)\bigr)\ \text{with corners }(2,4),(3,1),(5,0),(4,3).
\end{align*}
Each has area $5=d$ (Lemma~\ref{lem:det}), so their areas add to $20$, whereas the area of their union is $25/2$. This is the overlap phenomenon of Example~\ref{ex:overlap}: the cells are distinct admissible parallelograms, not a tiling, and the union is what the theory measures.

The complement $H_0\setminus U_0$ is $\interior T_1\sqcup\interior T_2$ with
\[
  T_1=\conv\bigl((0,5),(2,4),(5,5)\bigr),\qquad T_2=\conv\bigl((5,0),(4,3),(5,5)\bigr),
\]
of area $5/2$ each, so the area of the holes is $5=d$ and $35/2-5=25/2$ is the area of $U_0$, in agreement. The centroids $(7/3,14/3)$ and $(14/3,8/3)$ of the two triangles lie in no admissible cell and on no admissible segment (checked in exact arithmetic). By Theorem~\ref{thm:betti}, $\beta_1(U_0)=2$ and $\beta_2(U_0)=0$.

\subsection{\texorpdfstring{The occurrences of $x_1$: one more point, no holes}{The occurrences of x1: one more point, no holes}}

The number $x_1=5^{12}3^{10}/11$ has $q_1=6$, $q_2=5$, so the normalized box is $Q=(6,5)$. It has one more fibre point, $(6,2)$ (note $2\equiv2\cdot6\bmod5$), and so eight occurrences (Table~\ref{tab:example}). The fibre is the former one together with $(6,2)$, the hull is
\[
  H_1=\conv\bigl((0,5),(1,2),(3,1),(5,0),(6,2),(5,5)\bigr),\qquad\text{area }20.
\]
Here $u=1$, $v=0$, so we are in the one-sided horizontal range of Theorem~\ref{thm:cap}. The outer chain is
\[
  \Omega=\bigl[(5,0),(6,2),(5,5),(0,5)\bigr]\quad\text{with classes }R,R,C,T.
\]
The first boundary step is $M=(1,2)$, the unique record of $2i\bmod5$ for $i\le u=1$, and the Graver vector $g=(1,-3)$ satisfies $\cross(g,M)=5\ge xy=2$. The cells that cover $H_1\setminus H_0$, and in addition the old holes, are:
\begin{itemize}
\item the \emph{bottom cell} $P(B-g;g,M)$ with corners $(4,3),(5,0),(6,2),(5,5)$;
\item the \emph{L-cell} on the edge $(6,2)\to(5,5)$, with corners $(6,2),(5,5),(0,5),(1,2)$;
\item the two \emph{hole-filling cells} (the transposes of $R_1,R_2$ of Proposition~\ref{prop:fill} for $\Gamma_{5,3}$) with corners $(5,0),(4,3),(5,5),(6,2)$ and $(1,2),(6,2),(5,5),(0,5)$.
\end{itemize}
The first filling cell contains $T_2$ and the second contains $T_1$; for instance the vertex $(2,4)$ of $T_1$ lies inside the second one. In exact rational arithmetic the cells $C_1,\dots,C_4$ together with these four cover every point of the grid $\tfrac1{60}\Z^2\cap H_1$ (72\,301 points), as Theorem~\ref{thm:universal} predicts: $U_{5,2,(6,5)}=H_1$. In particular the realized union is convex and contractible, and the first Betti number has fallen from $2$ to $0$ by adding the single occurrence $(6,2)$.

\begin{remark}[The two values are not identified by inclusion]\label{rem:x0x1}
The exponent fibres satisfy $F_{x_0}\subset F_{x_1}$: every index $\kappa$ of an occurrence of $x_0$ is also an index of an occurrence of $x_1$. The occurrences themselves are different pairs. An occurrence determines its value, so no occurrence of $x_0$ is an occurrence of $x_1$. In Table~\ref{tab:example} the index $\kappa$ gives $(N,M)$ for $x_0$ and $(N,25M)$ for $x_1$. Hence the passage from the exceptional box to the next larger one is a statement about two fibres that share a lattice and a box-inclusion, and the identification of the two geometries is made at the level of the normalized lattice $\Gamma_{5,2}$ (Proposition~\ref{prop:transfer}), not at the level of occurrence pairs. This is the mechanism behind Question~\ref{q:global}: nothing in this paper says how the occurrences of different values are related to each other in a way that the geometry sees.
\end{remark}

\begin{remark}[The interior corner]\label{rem:ex-d9}
For the box $Q=(8,8)$ of $\Gamma_{5,2}$ (realized, for example, by $5^{16}3^{16}/11$) the point $C_0=(5,5)$ is an interior point of the hull, by the identity $(5,5)=\tfrac19(5,0)+\tfrac13(0,5)+\tfrac59(8,6)$ of Remark~\ref{rem:D9}, which the script verifies in exact arithmetic. The outer chain of this box therefore does not pass through $C_0$.
\end{remark}

%% -------------------------------------------------------------------
\section{Consequences for occurrence fibres}\label{sec:consequences}

This section collects the statements about actual occurrences that follow from the theorems of Parts~II and~III together with the exact arithmetic of Part~I. Each is a consequence in the strict sense; none introduces a new hypothesis.

\subsection{A threshold law for actual two-prime fibres}

Let $x=a/b$ be a reduced fraction whose numerator has exactly two prime factors $p_1,p_2$ coprime to $b$, with exponent bounds $q_i=\lfloor e_i/2\rfloor\ge1$, and assume $d\ge2$. Put $o_i=\ord_b(p_i)$ and $\lambda_i=o_i/d$. There are three regimes.

\begin{enumerate}
\item \emph{Thin} ($q_i<o_i$ for some $i$). The normalized box has a side below $d$, and our theorems say nothing.
\item \emph{Plateau} ($o_i\le q_i\le o_i+\lambda_i-1$ for $i=1,2$). The normalized box is $(d,d)$. The realized union has $\beta_1=d$ if $s\equiv-1\pmod d$, and $\beta_1=2$ otherwise; it is the fan graph, resp.\ the hull with two open triangles removed.
\item \emph{Convex} ($q_i\ge o_i$ for $i=1,2$ and not on the plateau). The realized union is the whole hull.
\end{enumerate}

This is Corollary~\ref{cor:actual}. It has the form of a threshold law: if one starts on the plateau and raises the exponents of the numerator (replacing $x$ by $xp_1^2$ or $xp_2^2$ repeatedly), the first Betti number of the realized union drops from $2$ (or $d$) to $0$ as soon as some normalized box side exceeds $d$, that is, when some $q_i$ reaches $o_i+\lambda_i$, and it stays $0$ afterwards (Corollary~\ref{cor:persistence}).

\begin{remark}[Sizes]
The plateau is not a small set of values. If $o_1=o_2=d$ (so $\lambda_i=1$) the plateau is the single bound pair $q=(d,d)$, that is, the exponent of $p_i$ in $a$ is $2d$ or $2d+1$. In general it is a rectangle of $\lambda_1\lambda_2$ bound pairs. The plateau is the first place at which both $q_i\ge o_i$ holds, and every value with both $q_i\ge o_i$ lies on it or beyond it.
\end{remark}

\subsection{Realizability and the family of normalized instances}

By Proposition~\ref{prop:realize} every pair $(d,s)$ occurs, with every box. Consequently Theorems~\ref{thm:betti} and~\ref{thm:universal} describe, for each modulus $d$ and each unit $s$, an actual set of values. In particular the two-hole phenomenon is not an artefact of the normalization: for every $d\ge3$ and every unit $s\not\equiv-1$ there is a rational number whose realized union of Graver cells has exactly two bounded complementary components, and the fan graph with $\beta_1=d$ occurs for every $d\ge2$ as the exceptional topology of the $s\equiv-1$ class. We repeat that Proposition~\ref{prop:realize} is an addition of this manuscript and that no result depends on it.

\subsection{Arithmetic meaning of the cells}\label{sec:cell-arith}

Let $w,w'\in F_x$ be two points, $\kappa=\kappa(w)$ and $\kappa'=\kappa(w')$ the corresponding indices. Then $\kappa'/\kappa=p_1^{w'_1-w_1}p_2^{w'_2-w_2}$. Because both indices are $\equiv1\pmod b$, the vector $g=w'-w$ lies in $\Lambda_b$, that is, $p_1^{g_1}p_2^{g_2}\equiv1\pmod b$ (with a negative exponent interpreted in $(\Z/b)^\times$). We distinguish two situations.

\begin{itemize}
\item If $g\ge0$ coordinatewise, then $\kappa\mid\kappa'$ and the quotient $\kappa'/\kappa=p_1^{g_1}p_2^{g_2}$ is itself $\equiv1\pmod b$. Such $g$ is Graver if and only if it is a \emph{minimal} nonnegative relation: no proper divisor of the quotient, other than $1$, is $\equiv1\pmod b$ and of the form $p_1^{h_1}p_2^{h_2}$ with $h\le g$.
\item If $g=(g_1,-g_2)$ with $g_1,g_2>0$ then the relation reads $\kappa'\,p_1^{g_1}=\kappa\,p_2^{g_2}$: two indices related by a common multiple. Graver means minimal for the sign pattern.
\end{itemize}

A cell $P(w;g,h)$ with all four corners in $F_x$ is therefore a ``commuting square'' of occurrences: it records two minimal relations $g$ and $h$ that can be applied in either order, starting from the occurrence $w$, without leaving the set of occurrences. The statement $U=\conv F$ says that every point of the convex hull of the occurrence fibre in exponent space lies in such a commuting square or on a minimal-relation segment. The arithmetic content of the stabilization theorem is thus a statement of local abundance of commuting squares once the box exceeds $(d,d)$, and its failure in the exceptional box is localized in two explicit triangles that have no commuting squares through them.

\begin{remark}[Not a statement about all pairs]
We do not claim that every pair of occurrences in the same fibre is connected by a chain of such squares in the sense of an arithmetic relation among $N$ and $M$. The set $U_Q$ lives in exponent space and records which convex combinations of exponent vectors lie in a cell. How these convex combinations are to be interpreted as arithmetic objects (they are not integers) is not addressed here.
\end{remark}

\subsection{What the theorems determine about the hull}

Since $U_Q=H_Q=\conv F_Q$ beyond the exceptional box, the realized union is a convex lattice polygon with vertices in $F_Q$, and so all its invariants are determined by the extreme points of the fibre. In particular:

\begin{enumerate}
\item the area of $U_Q$ is the area of the polygon, a half-integer because the vertices are integer points; measured in units of the covolume $d$ of $\Gamma$ it is $\operatorname{area}(H_Q)/d$;
\item the boundary of $U_Q$ is the polygonal line through the outer chain $\Omega$ and the lower-left chain, every edge of which is a Graver segment (Lemma~\ref{lem:reflection});
\item $\beta_1(U_Q)=0$ for every box $Q\ge(d,d)$, $Q\ne(d,d)$ (Corollary~\ref{cor:contractible}), so the exceptional first Betti numbers $2$ and $d$ are visible at exactly one box.
\end{enumerate}

In the exceptional box the deficiency $\operatorname{area}(H_0)-\operatorname{area}(U_0)$ equals $d$ in the two-hole classes (two triangles of area $d/2$) and $d^2/2$ (the whole hull) in the class $s\equiv-1$. The jump of the first Betti number and the area deficiency are thus both controlled by the same two triangles $T_1$, $T_2$ of Definition~\ref{def:T}.

\subsection{A structural summary}

Table~\ref{tab:summary} summarises the regimes for the normalized family. The first column gives the box, the second the modular class, the third the realized union and the fourth the theorem that proves it.

\begin{table}[ht]
\centering\small
\begin{tabular}{@{}llll@{}}
\toprule
Box & Class & $U_Q$ & Where proved\\
\midrule
$Q=(d,d)$ & $d\ge5$, $2\le s\le d-2$ & $H_0\setminus(\interior T_1\sqcup\interior T_2)$, $\beta_1=2$ & Thm~\ref{thm:exceptional}, \ref{thm:betti}\\
$Q=(d,d)$ & $s\equiv+1$, $d\ge3$ & $W\cup3$ segments, $\beta_1=2$ & Prop.~\ref{prop:special-plus}\\
$Q=(d,d)$ & $s\equiv-1$, $d\ge2$ & fan graph, $\beta_1=d$ & Prop.~\ref{prop:special-fan}\\
$Q\ge(d,d)$, $Q\ne(d,d)$ & all & $\conv F_Q$, contractible & Thm~\ref{thm:universal}\\
some $Q_i<d$ & all & not treated & ---\\
\bottomrule
\end{tabular}
\caption{The regimes of the realized union for the normalized family $\Gamma_{d,s}$.}\label{tab:summary}
\end{table}

%% -------------------------------------------------------------------
\section{Literature and context}\label{sec:lit}

\paragraph{Scope of this review and a disclosure.}
This review maps the neighbouring mathematics so that a reader can see where the present construction sits. It makes no priority or novelty claim, and no systematic search of the literature on lattice geometry was carried out. The bibliographic data of the works of Graver, Hemmecke, Richter-Gebert--Ziegler, Sturmfels, Schrijver, Ziegler, Beck--Robins and Karpenkov were checked against publisher, society, catalogue or open-access index pages; for these works \emph{only abstracts or catalogue pages were read, not the full texts}, and we have not verified any theorem number inside them. The data for Hatcher and for Bj\"orner--Tancer were taken from a reference list on an encyclopedia page about Alexander duality, and the data for Serre are given from the author's knowledge; they were not checked against a catalogue. Where a sentence below says what a notion is, it records what the cited source's title and abstract lead us to understand, and is not a quotation of a result.

\paragraph{Graver bases and conformal minimality.}
The set of conformally minimal elements of a lattice, taken with respect to the coordinatewise sign-compatible order $h\conf g$ of Definition~\ref{def:conformal}, is the Graver basis of the lattice, introduced by Graver in the study of the foundations of integer linear programming~\cite{Graver1975}. For a lattice of the form $\ker_\Z A$ it is finite. Its connection with toric ideals and Gr\"obner bases is treated in Sturmfels's lecture notes~\cite{Sturmfels1996}; Schrijver's monograph~\cite{Schrijver1986} is a general reference for the integer-programming background. Algorithms for computing Graver bases, among them the positive sum property, are given by Hemmecke~\cite{Hemmecke2003}. In our setting the lattice is two-dimensional, and Lemma~\ref{lem:std-graver} exhibits by hand six Graver vectors, $\pm(d,0)$, $\pm(0,d)$, $\pm(1,s)$, $\pm(1,s-d)$, $\pm(t,1)$, $\pm(d-t,-1)$; the full Graver basis may be larger (Appendix~\ref{app:graver}), and the constructions below use only the Graver property of the generators they employ. What is specific to this paper is not the Graver basis but its use to define an \emph{attachment rule} of parallelogram cells with corners in a finite box fibre.

\paragraph{Parallelograms, zonogons and zonotopal tilings.}
A parallelogram spanned by two vectors is the simplest zonotope, and a zonogon is a Minkowski sum of segments. The classical tilings of zonogons by parallelograms and their correspondence with oriented-matroid data are the subject of the Bohne--Dress theorem; see Richter-Gebert and Ziegler~\cite{RichterGebertZiegler1994} and Ziegler's text on polytopes~\cite{Ziegler1995}. Our zonogon $Z$ of Lemma~\ref{lem:zonogon} is tiled by the sweep of Graver steps of the lower-left hull chain, and the proof there is a direct induction on the number of steps, which includes the case of parallel steps; we have not used the Bohne--Dress theorem. The relation between our sweep tiling and the standard rhombic tilings is plausible but not examined here.

\paragraph{Lattice points in polytopes and the geometry of numbers.}
Enumeration of lattice points in polytopes, including Pick-type area formulas and Ehrhart theory, is covered by Beck and Robins~\cite{BeckRobins2015}. The fibre $F_Q$ is the set of nonzero lattice points of a congruence sublattice in a box, and the hull $H_Q$ is a lattice polygon for the lattice $\Gamma_{d,s}$; the arithmetic inequality of Lemma~\ref{lem:arith} bounds the determinants of triangles spanned by Pareto-minimal points of $\Gamma_{d,s}$ and is, in spirit, a statement of the geometry of numbers. We do not claim a precise relation to the classical theorems (Minkowski's convex body theorem, Dirichlet approximation).

\paragraph{Continued fractions and sails.}
The lower-left chain of Lemma~\ref{lem:chain0} and the record minima of $si\bmod d$ in Section~\ref{sec:pareto} are reminiscent of the sails of Klein polygons and of the continued-fraction expansion of $s/d$; see Karpenkov's book~\cite{Karpenkov2013} on the geometry of continued fractions. The chain of Lemma~\ref{lem:chain0}, with steps $(1,-a)$, \dots, $(c,-1)$ between the points $A$ and $B$, is the boundary of the convex hull of the nonzero lattice points of $\Gamma_{d,s}$ in the positive quadrant, i.e.\ the lower-left part of such a sail, restricted to the box. We have not established a precise dictionary, in particular between the Graver steps and the partial quotients of $s/d$, and we state it only as an analogy.

\paragraph{Homology of planar sets.}
For the topological counts of Section~\ref{sec:topology} we use Alexander duality for compact polyhedra in the plane, in the form of Hatcher~\cite[Theorem~3.44, \S3.3]{Hatcher2002} (the theorem number was confirmed by an independent verification review against an online text, not by us from the printed book). Bj\"orner and Tancer~\cite{BjornerTancer2009} give a short proof of \emph{combinatorial} Alexander duality for simplicial complexes; it is related background and is not used here as a substitute for the topological statement quoted from Hatcher. We state duality only for compact polyhedra in the plane and do not need the general case.

\paragraph{Occurrence arithmetic.}
The occurrence equation $x=M(N+1)^2/N$ is specific to the UNNS nest, and we have not found a direct antecedent in the literature. The reduction to $\kappa^2\mid a$, $\kappa\equiv1\pmod b$ (Proposition~\ref{prop:bij}) is elementary. Congruence lattices as kernels of homomorphisms to $(\Z/b)^\times$ are standard; for the cyclic structure of the unit groups and Dirichlet's theorem, used in Proposition~\ref{prop:realize}, see for example~\cite{Serre1973}.

\paragraph{What is studied here.}
The objects studied in this paper are the occurrence-coordinate correspondence for the fibre of a single rational value, the normalized modular lattices $\Gamma_{d,s}$ obtained from it, and the realized Graver-cell union on the punctured box fibre. The theorems concern this construction only. Comparisons with the topics above are offered as orientation: the lower-left chain is analogous to a sail, the cells are zonotopal, the generators are Graver vectors, and the exceptional holes are the first place where the cell construction fails to reach the convex hull. No further relation is asserted.

\paragraph{Scope note on the term ``arithmetic occurrence geometry''.}
The expression is used here as a defined term (Definition~\ref{def:aog}) for the arithmetic-geometric constructions studied in this work. It is not intended as a synonym for the established field of arithmetic geometry, nor does its use imply a claim of priority over every previous occurrence of similar terminology. We did not survey other fields for similar wording; we understand that phrases of the form ``occurrence geometry'' are used with unrelated meanings elsewhere (for instance for spatial patterns of species occurrence in ecology), and the present meaning is only the one fixed by Definition~\ref{def:aog}.

\subsection{Primary sources: the pinned GEO-I archive}\label{sec:lit-archive}

The mathematical development behind this paper is recorded in twelve archived research releases of the same research programme, GEO-I v001--v012, held in the repository of the UNNS project. All citations below are to the fixed commit
\begin{center}\geoicommit\end{center}
so that a reader can inspect the exact source versions used in preparing the manuscript, even if the main branch later changes. Each release is cited in the bibliography~\cite{GEOIv001,GEOIv002,GEOIv003,GEOIv004,GEOIv005,GEOIv006,GEOIv007,GEOIv008,GEOIv009,GEOIv010,GEOIv011,GEOIv012} as an \emph{archived research release}. These are primary sources for our own programme. They are not independently peer-reviewed publications, and they do not replace the external mathematical literature cited above for Graver bases, lattice geometry, zonogons and algebraic topology. Their role in this paper is given in Table~\ref{tab:geoi-roles}; the precise theorem-level citations are in Table~\ref{tab:geoi-theorems} and in the registry of Appendix~\ref{app:registry}.

\begin{table}[ht]
\centering\small
\begin{tabular}{@{}lp{12.2cm}@{}}
\toprule Release & Role in this paper\\\midrule
v001~\cite{GEOIv001} & Arithmetic occurrence fibres and initial geometric constructions\\
v002~\cite{GEOIv002} & Arithmetic occurrence fibres and initial geometric constructions\\
v003~\cite{GEOIv003} & Abstract versus embedded geometry; finite obstruction theory\\
v004~\cite{GEOIv004} & Abstract versus embedded geometry; planar normalization (Prop.~S1)\\
v005~\cite{GEOIv005} & Special-class stabilization G5-C1 and exceptional topology G5-C2\\
v006~\cite{GEOIv006} & Preliminary modular geometry and computational investigations (context)\\
v007~\cite{GEOIv007} & Large-box theorem G7-T1; exceptional-hole lower bound G7-T3\\
v008~\cite{GEOIv008} & First-enlargement hole filling G8-T1; earlier coverage reduction G8-R1\\
v009~\cite{GEOIv009} & Intermediate geometric results and historical proof development (context)\\
v010~\cite{GEOIv010} & Exact exceptional complement G10-T1; old-hull coverage G10-R1; hull chain G10-L1, G10-Z\\
v011~\cite{GEOIv011} & Intermediate planar results (G11-A etc.) and conjectures, resolved or corrected in v012\\
v012~\cite{GEOIv012} & Pareto inequality G12-T1, boundary Graver theorem G12-T2, cap cover G12-T3--T7, universal stabilization G12-T8\\
\bottomrule
\end{tabular}
\caption{The twelve GEO-I releases and their roles. The principal mathematical sources are v005, v007, v008, v010 and v012; v006, v009 and v011 are historical context.}\label{tab:geoi-roles}
\end{table}

\begin{table}[ht]
\centering\small
\begin{tabular}{@{}p{3.3cm}p{5.4cm}p{5.8cm}@{}}
\toprule Identifier & Release and document & Used in this paper for\\\midrule
G5-C1, G5-C2 & v005~\cite{GEOIv005}, \texttt{STABILIZATION.md} & Theorem~\ref{thm:L}; Proposition~\ref{prop:special-fan}; $\beta_1=2$ for $s=+1$\\
G7-T1 & v007~\cite{GEOIv007}, \texttt{COVERAGE.md} & Theorem~\ref{thm:large}\\
G7-T3 & v007~\cite{GEOIv007}, \texttt{HOLLOW\_TOPOLOGY.md} & Proposition~\ref{prop:hole-exclusion}\\
G8-T1, G8-R1 & v008~\cite{GEOIv008}, \texttt{HOLLOW.md}, \texttt{SMALL\_BOX.md} & Proposition~\ref{prop:fill}; the reduction behind Corollary~\ref{cor:old}\\
G10-T1, G10-T1H, G10-R1, G10-Z, G10-L1 & v010~\cite{GEOIv010}, \texttt{MATHEMATICS.md}, \texttt{HOLLOW.md} & Theorem~\ref{thm:exceptional}; Theorem~\ref{thm:betti}; Corollary~\ref{cor:old}; Lemmas~\ref{lem:chain0}, \ref{lem:zonogon}\\
G12-T1 -- G12-T8, G12-L1, L4, L5 & v012~\cite{GEOIv012}, \texttt{MATHEMATICS.md}, \texttt{MIXED\_CAP.md}, \texttt{GRAVER\_STRIP.md} & Lemmas~\ref{lem:arith}, \ref{lem:reflection}; Theorem~\ref{thm:cap}; Theorem~\ref{thm:universal}\\
\bottomrule
\end{tabular}
\caption{Theorem-level source citations (identifier, release, document). The exact-complement statement of Proposition~\ref{prop:special-plus} is a contribution of this manuscript and is not attributed to v005 (Remark~\ref{rem:s1-source}).}\label{tab:geoi-theorems}
\end{table}
